\paragrafos{Άθροισμα γωνιών τριγώνου}{Β.3.2, «Άθροισμα γωνιών τριγώνου – Ιδιότητες ισοσκελούς τριγώνου»}

\begin{anagke}
Ο Βασίλης κόβει από χαρτί ένα τρίγωνο και σκίζει τις τρεις γωνίες του. Τις βάζει δίπλα δίπλα, με τις κορυφές στο ίδιο σημείο, και οι τρεις μυτούλες σχηματίζουν μια ευθεία γραμμή. «Ευθεία γωνία!», φωνάζει. Η Αγγελική κάνει το ίδιο με ένα άλλο τρίγωνο, πιο λεπτό και μακρύ: πάλι ευθεία. Ο Νίκος το κάνει με αμβλυγώνιο: ξανά ευθεία.

Η Πολυξένη ρωτάει: «Ισχύει πάντα; Και πώς το ξέρουμε χωρίς να κόψουμε όλα τα τρίγωνα του κόσμου;». Η απάντηση είναι ότι το αποδεικνύουμε με όσα ξέρουμε για τις παράλληλες ευθείες. Η Ρήνη προσθέτει: «Και τα ισοσκελή; Αν διπλώσω ένα ισοσκελές στη μέση, οι δύο μισές πέφτουν η μία πάνω στην άλλη!».

Σε αυτή την παράγραφο μαθαίνουμε ότι το άθροισμα των γωνιών κάθε τριγώνου είναι $180^\circ$, τι ισχύει για την εξωτερική γωνία, και τις ιδιότητες του ισοσκελούς και του ισόπλευρου τριγώνου.
\end{anagke}

\begin{lexi}{εξωτερική γωνία, βάση}{ονόματα}
Η \lexiword{εξωτερική γωνία} είναι αυτή που σχηματίζεται από μια πλευρά και την προέκταση της διπλανής. Η \lexiword{βάση} του ισοσκελούς είναι η πλευρά που δεν είναι από τις δύο ίσες, και οι \lexiword{προσκείμενες γωνίες} της είναι οι δύο γωνίες στα άκρα της.
\end{lexi}

\begin{metafrasi}
Οι τρεις γωνίες φτιάχνουν ευθεία γωνία. & $\hat{\mathrm{Α}}+\hat{\mathrm{Β}}+\hat{\Gamma}=180^\circ$\\
Οι δύο οξείες του ορθογωνίου τριγώνου. & συμπληρωματικές: $\hat{\mathrm{Β}}+\hat{\Gamma}=90^\circ$\\
Η εξωτερική γωνία στο $\Gamma$. & $\hat{\varphi}=\hat{\mathrm{Α}}+\hat{\mathrm{Β}}$\\
Οι γωνίες στη βάση του ισοσκελούς. & $\hat{\mathrm{Β}}=\hat{\Gamma}$\\
Όλες οι γωνίες του ισόπλευρου. & $60^\circ$\\
\end{metafrasi}

\tmima{Η θεωρία}

\begin{kanonas}{άθροισμα γωνιών τριγώνου}
Σε κάθε τρίγωνο $\mathrm{ΑΒΓ}$: $\hat{\mathrm{Α}}+\hat{\mathrm{Β}}+\hat{\Gamma}=180^\circ$.
\end{kanonas}

\textbf{Γιατί ισχύει.} Φέρνω από το $\mathrm{Α}$ την ευθεία $x\mathrm{Α}y$ παράλληλη προς τη $\mathrm{ΒΓ}$. Η $x\hat{\mathrm{Α}}\mathrm{Β}$ είναι ίση με την $\hat{\mathrm{Β}}$ (εντός εναλλάξ) και η $y\hat{\mathrm{Α}}\Gamma$ ίση με την $\hat{\Gamma}$ (εντός εναλλάξ). Οι τρεις γωνίες γύρω από το $\mathrm{Α}$ σχηματίζουν ευθεία γωνία, άρα $\hat{\mathrm{Β}}+\hat{\mathrm{Α}}+\hat{\Gamma}=180^\circ$.

\begin{center}
\begin{tikzpicture}[scale=0.9]
\draw[thick] (-0.4,0)--(3.8,0) ;
\node[below left] at (0,0) {$\mathrm{Β}$}; \node[below right] at (3.4,0) {$\Gamma$};
\draw[thick] (0,0)--(1.4,1.8)--(3.4,0)--cycle;
\draw[thick,gblue] (-0.6,1.8)--(3.6,1.8);
\node[above] at (1.4,1.8) {$\mathrm{Α}$};
\node at (0.55,0.2) {\small $\omega$}; \node at (2.9,0.2) {\small $\theta$};
\node at (0.7,1.6) {\small $\omega$}; \node at (2.1,1.6) {\small $\theta$};
\node at (1.4,1.45) {\small $\hat{\mathrm{Α}}$};
\end{tikzpicture}
\end{center}

\begin{kanonas}{συνέπειες}
\begin{itemize}[leftmargin=1.2em,itemsep=0.3ex,topsep=2pt]
\item Στο \textbf{ορθογώνιο} τρίγωνο οι δύο οξείες γωνίες είναι \textbf{συμπληρωματικές}: $\hat{\mathrm{Β}}+\hat{\Gamma}=90^\circ$.
\item Η \textbf{εξωτερική γωνία} τριγώνου ισούται με το άθροισμα των δύο \emph{απέναντι} εσωτερικών: $\hat{\varphi}=\hat{\mathrm{Α}}+\hat{\mathrm{Β}}$ (γιατί $\hat{\varphi}=180^\circ-\hat{\Gamma}$).
\item Στο \textbf{ισόπλευρο} τρίγωνο κάθε γωνία είναι $180^\circ:3=60^\circ$.
\end{itemize}
\end{kanonas}

\begin{kanonas}{ισοσκελές τρίγωνο}
Σε ισοσκελές τρίγωνο ($\mathrm{ΑΒ}=\mathrm{ΑΓ}$, κορυφή $\mathrm{Α}$, βάση $\mathrm{ΒΓ}$):
\begin{itemize}[leftmargin=1.2em,itemsep=0.3ex,topsep=2pt]
\item η ευθεία της διαμέσου προς τη βάση είναι \textbf{άξονας συμμετρίας},
\item η διάμεσος προς τη βάση είναι και \textbf{ύψος} και \textbf{διχοτόμος},
\item οι \textbf{προσκείμενες στη βάση γωνίες είναι ίσες}: $\hat{\mathrm{Β}}=\hat{\Gamma}$.
\end{itemize}
Στο \textbf{ισόπλευρο}: οι ευθείες των διαμέσων είναι άξονες συμμετρίας, κάθε διάμεσος είναι ύψος και διχοτόμος, όλες οι πλευρές και γωνίες είναι ίσες.
\end{kanonas}

\begin{prosoxi}
\begin{itemize}[leftmargin=1.2em,itemsep=0.3ex,topsep=2pt]
\item Στο ισοσκελές, η ιδιότητα «διάμεσος = ύψος = διχοτόμος» ισχύει για το τμήμα προς τη \emph{βάση}, όχι για τα άλλα.
\item Στο ισοσκελές ορθογώνιο τρίγωνο οι δύο οξείες είναι $45^\circ$ η καθεμία.
\item Η εξωτερική γωνία είναι ίση με το άθροισμα των \emph{δύο άλλων} γωνιών (όχι της διπλανής της).
\end{itemize}
\end{prosoxi}

\tmima{Λυμένα παραδείγματα}

\paradeigma{Η τρίτη γωνία}
\tixero Τρίγωνο με γωνίες $\hat{\mathrm{Α}}=47^\circ$ και $\hat{\mathrm{Β}}=68^\circ$.\par
\tithelo Η $\hat{\Gamma}$.\par
\possindeo $\hat{\Gamma}=180^\circ-(47^\circ+68^\circ)=180^\circ-115^\circ=65^\circ$.\par
\apantisi $65^\circ$ (οξυγώνιο τρίγωνο).\par
\elegxos $47+68+65=180$.

\paradeigma{Η εξωτερική γωνία}
\tixero Τρίγωνο με $\hat{\mathrm{Α}}=50^\circ$, $\hat{\mathrm{Β}}=60^\circ$.\par
\tithelo Η εξωτερική γωνία στο $\Gamma$.\par
\possindeo $\hat{\varphi}=\hat{\mathrm{Α}}+\hat{\mathrm{Β}}=110^\circ$. Ή: $\hat{\Gamma}=70^\circ$ και $\hat{\varphi}=180^\circ-70^\circ=110^\circ$.\par
\apantisi $110^\circ$.\par
\elegxos Οι δύο τρόποι δίνουν το ίδιο.

\paradeigma{Ισοσκελές με γωνία κορυφής}
\tixero Ισοσκελές τρίγωνο με γωνία κορυφής $\hat{\mathrm{Α}}=40^\circ$.\par
\tithelo Οι γωνίες της βάσης.\par
\possindeo $\hat{\mathrm{Β}}=\hat{\Gamma}$ και $\hat{\mathrm{Β}}+\hat{\Gamma}=180^\circ-40^\circ=140^\circ$. Άρα $\hat{\mathrm{Β}}=\hat{\Gamma}=70^\circ$.\par
\apantisi $70^\circ$ και $70^\circ$.\par
\elegxos $40+70+70=180$.

\paradeigma{Ορθογώνιο και ισοσκελές}
\tixero Ένα τρίγωνο είναι ορθογώνιο και ισοσκελές.\par
\tithelo Οι γωνίες του.\par
\possindeo Η ορθή είναι $90^\circ$. Οι άλλες δύο είναι ίσες και με άθροισμα $90^\circ$, άρα $45^\circ$ η καθεμία.\par
\apantisi $90^\circ,45^\circ,45^\circ$.\par
\elegxos $90+45+45=180$.

\tmima{Ασκήσεις}

\zoni{Το κατάλαβα}{γωνίες τριγώνων}

\ex Βρες την τρίτη γωνία τριγώνου με γωνίες: α) $35^\circ,\,80^\circ$ β) $90^\circ,\,28^\circ$ γ) $112^\circ,\,41^\circ$.

\ex Ισοσκελές τρίγωνο έχει γωνία κορυφής $50^\circ$. Βρες τις γωνίες της βάσης.

\ex Ισοσκελές τρίγωνο έχει γωνία στη βάση $72^\circ$. Βρες τη γωνία κορυφής.

\ex Πόσο είναι κάθε γωνία ισοπλεύρου τριγώνου; Πόσο η εξωτερική του γωνία;

\ex Οι οξείες ενός ορθογωνίου τριγώνου είναι $x$ και $2x$. Βρες τες.

\ex Σε τρίγωνο η εξωτερική γωνία στο $\Gamma$ είναι $125^\circ$ και $\hat{\mathrm{Α}}=48^\circ$. Βρες τη $\hat{\mathrm{Β}}$ και τη $\hat{\Gamma}$.

\ex Σωστό ή λάθος; α) Τρίγωνο μπορεί να έχει δύο ορθές γωνίες. \quad β) Οι γωνίες βάσης ισοσκελούς είναι ίσες. \quad γ) Σε ισόπλευρο τρίγωνο όλες οι γωνίες είναι $60^\circ$. \quad δ) Οι οξείες γωνίες ορθογωνίου είναι παραπληρωματικές.

\zoni{Το χρησιμοποιώ}{στέγες και σήματα}

\ex \textbf{Η στέγη.} Μια στέγη έχει σχήμα ισοσκελούς τριγώνου με γωνία κορυφής $100^\circ$. Πόση είναι η γωνία που σχηματίζει κάθε πλευρά με το οριζόντιο δοκάρι;

\ex \textbf{Το σήμα.} Ένα τριγωνικό σήμα έχει ισόπλευρο σχήμα. Πόσο είναι το άθροισμα των τριών γωνιών του;

\ex \textbf{Η σκάλα.} Μια σκάλα ακουμπά στον τοίχο και σχηματίζει με αυτόν γωνία $25^\circ$. Ο τοίχος είναι κάθετος στο έδαφος. Πόση γωνία σχηματίζει η σκάλα με το έδαφος;

\ex \textbf{Ο χαρταετός.} Ένα τριγωνικό κομμάτι ενός χαρταετού έχει γωνίες $55^\circ$ και $75^\circ$. Πόση είναι η τρίτη γωνία;

\ex \textbf{Η γωνιά.} Ένας ξύλινος γνώμονας έχει σχήμα ισοσκελούς ορθογωνίου τριγώνου. Πόσες μοίρες είναι οι δύο ίσες γωνίες του;

\zoni{Το σκέφτομαι}{γιατί συμβαίνει}

\ex Εξήγησε με παράλληλη από την κορυφή γιατί το άθροισμα των γωνιών ενός τριγώνου είναι $180^\circ$.

\ex Εξήγησε γιατί τρίγωνο δεν μπορεί να έχει δύο αμβλείες γωνίες.

\ex Αν σε ισοσκελές τρίγωνο μια γωνία είναι $100^\circ$, ποια θέση έχει (κορυφής ή βάσης); Βρες τις άλλες δύο.

\ex Στο ισοσκελές τρίγωνο $\mathrm{ΑΒΓ}$ ($\mathrm{ΑΒ}=\mathrm{ΑΓ}$) η $\mathrm{ΑΔ}$ είναι διάμεσος προς τη $\mathrm{ΒΓ}$ και $\mathrm{ΒΓ}=12$ cm. Πόσο είναι το $\mathrm{ΒΔ}$ και πόση η γωνία $\mathrm{ΑΔΒ}$;

\zoni{Για το σπίτι}{γωνίες με το χέρι}

\ex Κόψε ένα τρίγωνο από χαρτί, σκίσε τις γωνίες του και τοποθέτησέ τες δίπλα δίπλα. Τι σχηματίζουν;

\ex Δίπλωσε ένα ισοσκελές τρίγωνο στη μέση και δες ποιες γωνίες πέφτουν η μία πάνω στην άλλη.
